% GENERATED FILE. Edit canonical lessons, then run npm run generate:books.
% canonical-source-hash: fnv1a64:063cb25b272b45fc
% canonical-lessons: FA-C434-pelle-barqi, FA-C434-masir, FA-C434-keraye, FA-C434-ja-gereftan, FA-C434-ja-mandan

\chapter{Underground and buses}
\label{ch:fa-underground-and-buses}
\hlchaptermodality{434}

\begin{chapteropening}
\textbf{By the end of this chapter:} \emph{I can get around by metro and bus: the escalator, the route, the fare, getting a seat and missing the bus.}

Complete the last lesson of chapter 434: I can get around by metro and bus: the escalator, the route, the fare, getting a seat and missing the bus.
\end{chapteropening}

\section[\texorpdfstring{\fa{پله} \fa{برقی} (pelle-barqi)}{پله برقی (pelle-barqi)}]{\fa{پله} \fa{برقی} (pelle-barqi) — an escalator}

\label{lesson:FA-C434-pelle-barqi}



\begin{quote}
Before the new one: say the Persian for to meet, then the Persian for to hurry.
\end{quote}

\subsection*{You'll want to know: \fa{پله} \fa{برقی}}
\textbf{\fa{پله} \fa{برقی}} (\emph{pelle-barqi}) — \textquotedblleft{}an escalator\textquotedblright{}.

\begin{quote}
\fa{در} \fa{ایستگاه} \fa{مترو} \fa{پله} \fa{برقی} \fa{خراب} \fa{بود}.
\end{quote}

\begin{quote}
\emph{Dar istgâh-e metro pelle-barqi kharâb bud.}
\end{quote}

\begin{quote}
The escalator at the metro station was out of order.
\end{quote}

Literally 'electric stairs': \emph{pelle} 'step' + \emph{barqi} 'electric'.

\begin{cousinweb}
\emph{pelle-barqi} is built from \emph{pelle} 'step, stair', whose origin is not settled, and \emph{barqi} 'electric', from \emph{barq} 'electricity', borrowed from Arabic \emph{barq} 'lightning, flashing' (EtymDB).
\end{cousinweb}

\begin{grammarlens}[title={What you've built}]
One more word for telling, explaining and describing at B1.
\end{grammarlens}

\subsection*{Your turn}
\begin{itemize}
\raggedright
  \item \emph{Say it:} \emph{pelle-barqi}
  \item \emph{Say it:} \emph{pelle-barqi}, once more
  \item \emph{Say it:} a sentence of your own with \emph{pelle-barqi}
  \item \emph{Recall:} say the Persian for to meet, then the Persian for to hurry, then say \emph{pelle-barqi} again
\end{itemize}

\begin{tcolorbox}[breakable,colback=teal!4,colframe=teal!35!black,title={Before you move on}]
What does \fa{پله} \fa{برقی} mean? (\textquotedblleft{}An escalator\textquotedblright{}.) Say it once more.
\end{tcolorbox}
\section[\texorpdfstring{\fa{مسیر} (masir)}{مسیر (masir)}]{\fa{مسیر} (masir) — a route, a way}

\label{lesson:FA-C434-masir}



\begin{quote}
Before the new one: say the Persian for to hurry, then the Persian for an escalator.
\end{quote}

\subsection*{You'll want to know: \fa{مسیر}}
\textbf{\fa{مسیر}} (\emph{masir}) — \textquotedblleft{}a route, a way\textquotedblright{}.

\begin{quote}
\fa{مسیر} \fa{این} \fa{اتوبوس} \fa{خیلی} \fa{طولانی} \fa{است}.
\end{quote}

\begin{quote}
\emph{Masir-e in otobus kheyli tulâni ast.}
\end{quote}

\begin{quote}
This bus's route is very long.
\end{quote}

\emph{Dar masir} — on the way, en route.

\begin{cousinweb}
\emph{masir} was borrowed from Arabic and is usually traced to Arabic \emph{masīr} 'going, the way of a journey'. EtymDB does not list it.
\end{cousinweb}

\begin{grammarlens}[title={What you've built}]
One more word for telling, explaining and describing at B1.
\end{grammarlens}

\subsection*{Your turn}
\begin{itemize}
\raggedright
  \item \emph{Say it:} \emph{masir}
  \item \emph{Say it:} \emph{masir}, once more
  \item \emph{Say it:} a sentence of your own with \emph{masir}
  \item \emph{Recall:} say the Persian for to hurry, then the Persian for an escalator, then say \emph{masir} again
\end{itemize}

\begin{tcolorbox}[breakable,colback=teal!4,colframe=teal!35!black,title={Before you move on}]
What does \fa{مسیر} mean? (\textquotedblleft{}A route, a way\textquotedblright{}.) Say it once more.
\end{tcolorbox}
\section[\texorpdfstring{\fa{کرایه} (kerâye)}{کرایه (kerâye)}]{\fa{کرایه} (kerâye) — a fare; a rent, a hire charge}

\label{lesson:FA-C434-keraye}



\begin{quote}
Before the new one: say the Persian for an escalator, then the Persian for a route.
\end{quote}

\subsection*{You'll want to know: \fa{کرایه}}
\textbf{\fa{کرایه}} (\emph{kerâye}) — \textquotedblleft{}a fare; a rent, a hire charge\textquotedblright{}.

\begin{quote}
\fa{کرایه} \fa{تاکسی} \fa{را} \fa{پرداختیم}.
\end{quote}

\begin{quote}
\emph{Kerâye-ye tâksi râ pardâkhtim.}
\end{quote}

\begin{quote}
We paid the taxi fare.
\end{quote}

\emph{Kerâye kardan} — to hire (a car, a hall).

\begin{cousinweb}
\emph{kerâye} was borrowed from Arabic \emph{kirāʾ} 'hire, wage' (EtymDB).
\end{cousinweb}

\begin{grammarlens}[title={What you've built}]
One more word for telling, explaining and describing at B1.
\end{grammarlens}

\subsection*{Your turn}
\begin{itemize}
\raggedright
  \item \emph{Say it:} \emph{kerâye}
  \item \emph{Say it:} \emph{kerâye}, once more
  \item \emph{Say it:} a sentence of your own with \emph{kerâye}
  \item \emph{Recall:} say the Persian for an escalator, then the Persian for a route, then say \emph{kerâye} again
\end{itemize}

\begin{tcolorbox}[breakable,colback=teal!4,colframe=teal!35!black,title={Before you move on}]
What does \fa{کرایه} mean? (\textquotedblleft{}A fare; a rent, a hire charge\textquotedblright{}.) Say it once more.
\end{tcolorbox}
\section[\texorpdfstring{\fa{جا} \fa{گرفتن} (jâ gereftan)}{جا گرفتن (jâ gereftan)}]{\fa{جا} \fa{گرفتن} (jâ gereftan) — to get a seat; to take up room}

\label{lesson:FA-C434-ja-gereftan}



\begin{quote}
Before the new one: say the Persian for a route, then the Persian for a fare.
\end{quote}

\subsection*{You'll want to know: \fa{جا} \fa{گرفتن}}
\textbf{\fa{جا} \fa{گرفتن}} (\emph{jâ gereftan}) — \textquotedblleft{}to get a seat; to take up room\textquotedblright{}.

\begin{quote}
\fa{در} \fa{مترو} \fa{جا} \fa{گرفتم} \fa{و} \fa{نشستم}.
\end{quote}

\begin{quote}
\emph{Dar metro jâ gereftam va neshastam.}
\end{quote}

\begin{quote}
In the metro I got a seat and sat down.
\end{quote}

\emph{jâ} 'place' + \emph{gereftan} 'to take'. \emph{Barâ-ye kasi jâ gereftan} — to save someone a seat.

\begin{cousinweb}
\emph{jâ gereftan} is built from \emph{jâ} 'place', inherited from Middle Persian \emph{gyāg} (EtymDB), and the native verb \emph{gereftan} 'to take', usually traced to Old Persian \emph{grab-} 'to seize'.
\end{cousinweb}

\begin{grammarlens}[title={What you've built}]
One more word for telling, explaining and describing at B1.
\end{grammarlens}

\subsection*{Your turn}
\begin{itemize}
\raggedright
  \item \emph{Say it:} \emph{jâ gereftan}
  \item \emph{Say it:} \emph{jâ gereftan}, once more
  \item \emph{Say it:} a sentence of your own with \emph{jâ gereftan}
  \item \emph{Recall:} say the Persian for a route, then the Persian for a fare, then say \emph{jâ gereftan} again
\end{itemize}

\begin{tcolorbox}[breakable,colback=teal!4,colframe=teal!35!black,title={Before you move on}]
What does \fa{جا} \fa{گرفتن} mean? (\textquotedblleft{}To get a seat; to take up room\textquotedblright{}.) Say it once more.
\end{tcolorbox}
\section[\texorpdfstring{\fa{جا} \fa{ماندن} (jâ mândan)}{جا ماندن (jâ mândan)}]{\fa{جا} \fa{ماندن} (jâ mândan) — to miss (a bus, a train); to be left behind}

\label{lesson:FA-C434-ja-mandan}



\begin{quote}
Before the new one: say the Persian for a fare, then the Persian for to get a seat.
\end{quote}

\subsection*{You'll want to know: \fa{جا} \fa{ماندن}}
\textbf{\fa{جا} \fa{ماندن}} (\emph{jâ mândan}) — \textquotedblleft{}to miss (a bus, a train); to be left behind\textquotedblright{}.

\begin{quote}
\fa{دیر} \fa{رسیدم} \fa{و} \fa{از} \fa{اتوبوس} \fa{جا} \fa{ماندم}.
\end{quote}

\begin{quote}
\emph{Dir residam va az otobus jâ mândam.}
\end{quote}

\begin{quote}
I arrived late and missed the bus.
\end{quote}

\emph{jâ} 'place' + \emph{mândan} 'to stay'. Things can be left behind too: \emph{kifam dar tâksi jâ mând}.

\begin{cousinweb}
\emph{jâ mândan} is built from \emph{jâ} 'place', inherited from Middle Persian \emph{gyāg} (EtymDB), and \emph{mândan} 'to remain', which EtymDB traces to Middle Persian: to stay behind in one's place.
\end{cousinweb}

\begin{grammarlens}[title={What you've built}]
One more word for telling, explaining and describing at B1.
\end{grammarlens}

\subsection*{Your turn}
\begin{itemize}
\raggedright
  \item \emph{Say it:} \emph{jâ mândan}
  \item \emph{Say it:} \emph{jâ mândan}, once more
  \item \emph{Say it:} a sentence of your own with \emph{jâ mândan}
  \item \emph{Recall:} say the Persian for a fare, then the Persian for to get a seat, then say \emph{jâ mândan} again
\end{itemize}

\begin{tcolorbox}[breakable,colback=teal!4,colframe=teal!35!black,title={Before you move on}]
What does \fa{جا} \fa{ماندن} mean? (\textquotedblleft{}To miss (a bus, a train); to be left behind\textquotedblright{}.) Say it once more.
\end{tcolorbox}
